\section{Confidential I/O and Network Data-Path Defense}

Confidential computing protects data in use, but real workloads rarely remain
inside CPU registers and local DRAM. They use NICs, DPUs, SmartNICs, storage
controllers, GPUs, NPUs, CXL-attached memory, RDMA paths, and virtual switches.
These components move data through DMA engines, device queues, interrupts,
MMIO windows, link layers, and sometimes device-local runtimes. A CPU TEE or
confidential VM therefore provides only a partial boundary unless the I/O path
is integrated into the same trust model~\cite{riscv_cove_io_2026,arm_smmu_spec,sok-tee,zhu2020hetee,dhar2024cloudscale,zhou2024snic,giantsidi2025tnic}.

This section treats confidential I/O and network defense as a separate survey
lane. It does not cover generic network security. Instead, it focuses on
mechanisms that let a confidential workload safely use a device, network path,
remote memory tier, or accelerator without handing plaintext or control
authority back to an untrusted host.

\subsection{Threat Model and Data-Path Surfaces}

\subsubsection{Why Device Paths Break CPU-Only Isolation}

The basic device threat is simple: DMA-capable devices can access memory
without executing CPU load and store instructions. If the platform protects a
confidential VM only through CPU page tables, PMP, GPT/GPC, or monitor-managed
state, a device with insufficiently constrained DMA authority can still reach
protected pages~\cite{arm_smmu_spec,riscv_iommu_2023,riscv_iopmp_2026}. The
same concern applies to MMIO registers and queue metadata. A host that cannot
read a TVM's memory directly may still observe descriptors, packet buffers,
completion queues, or device status if those structures are not assigned and
protected as part of the confidential boundary.

Interrupts create a second surface. Device completion and error events must be
delivered to the protected workload without letting the host spoof, suppress, or
misroute security-critical signals. Work such as Devlore makes this explicit
for confidential VMs: interrupt ownership and delivery are part of device
ownership, not a secondary OS scheduling detail~\cite{bertschi2026devlore}.
For RISC-V, AIA and CoVE-IO provide the architectural vocabulary for trusted MSI
and interrupt handling, while Arm platforms rely on SMMU and interrupt
controller integration in the broader CCA system model~\cite{riscv_aia_2023,riscv_cove_io_2026,arm_smmu_spec}.

\subsubsection{Network and Offload-Specific Risks}

Network devices add their own complications. A confidential workload may use a
virtual NIC, SR-IOV virtual function, SmartNIC, DPU, programmable switch path,
or RDMA queue pair. In each case, the sensitive state is not only application
payload. It includes descriptors, keys, queue ownership, completion events,
flow steering rules, packet headers, and device firmware state. A secure
offload design must therefore answer four questions: which device or function
is the endpoint, what code and configuration is running there, which memory can
it access, and how packets or DMA buffers are protected while crossing the
fabric.

FOLIO illustrates one possible tradeoff for confidential VM networking by
aiming for high-performance networks without requiring trusted I/O
devices~\cite{li2024folio}. This is useful because it separates two design
paths. One path tries to make the device part of the trusted computing base
through assignment, measurement, and secure transport. The other path tries to
avoid trusting the device by keeping plaintext and sensitive control inside the
confidential VM while still obtaining high throughput. Both directions belong
in this survey because they address the same confidential-computing bottleneck:
network performance should not require silently expanding the trusted boundary
to an opaque device stack.

Network endpoint identity is a separate question from device identity. Even if
a NIC or DPU can be identified through SPDM-like mechanisms, a tenant still
needs to know whether the application channel carrying secrets terminates
inside the intended TEE software. TLS+RA addresses this endpoint problem by
binding remote attestation evidence to TLS channel establishment, so that a
secure channel is not merely adjacent to a TEE quote but cryptographically tied
to the attested endpoint~\cite{weinhold2025tlsra}. This matters for
confidential vNIC, proxy, gateway, and DPU-hosted services because TLS
termination can otherwise become a place where plaintext exits the confidential
boundary.

SmartNIC and NIC-local trust work pushes the same question lower in the stack.
\textbf{S-NIC} shows that a multi-tenant SmartNIC needs hardware isolation for
on-NIC RAM, cache, accelerators, DMA, and bus bandwidth before offloaded network
functions can be treated as tenant-confidential components~\cite{zhou2024snic}.
\textbf{TNIC} takes a different route by placing a minimal root of trust at the
NIC, providing transferable authentication and non-equivocation for trustworthy
distributed systems~\cite{giantsidi2025tnic}. These works are not substitutes
for VM memory protection, SPDM, or TDISP. They clarify what a NIC or DPU must
provide if it becomes a real endpoint of a confidential data path.

\subsection{Protocol Layer: Identity, Session Security, and Device State}

\subsubsection{SPDM as Device Identity and Measurement Substrate}

The \textbf{Security Protocol and Data Model (SPDM)} provides a protocol for
device authentication, measurement retrieval, version and capability
negotiation, and session establishment~\cite{dmtf_spdm_2025,dmtf_spdm_arch_2024}.
In confidential-computing systems, SPDM is valuable because the host's view of a
device is not enough. A tenant or trusted manager needs evidence about the
device identity and state before allowing that device to participate in a
protected data path.

SPDM secured messages add confidentiality and integrity for protocol exchanges
after session establishment~\cite{dmtf_spdm_secured_messages_2025}. SPDM over
TCP extends the transport setting so that SPDM can be used beyond only local
bus-oriented interactions~\cite{dmtf_spdm_tcp_2024}. These specifications do
not themselves define a complete confidential I/O architecture. Their role is
to supply the identity, measurement, and secure-channel substrate that higher
level mechanisms such as TDISP, CoVE-IO, and trusted I/O designs can rely on.

\subsubsection{TDISP and Trusted Device Interfaces}

\textbf{TDISP} addresses a more specific lifecycle problem: how a PCIe device
interface can be placed into a TEE-aware trust state and assigned to a
confidential workload~\cite{pcisig_tdisp_2022}. Device identity alone is not
enough. The platform also needs a way to bind a particular device interface,
function, or trusted device interface to the workload that will use it. TDISP
therefore belongs at the device-interface lifecycle layer, whereas SPDM belongs
at the identity and secure-session layer.

The RISC-V CoVE-IO draft uses this kind of structure explicitly. It refers to
trusted device interfaces, device managers, security managers, SPDM, TDISP,
PCIe IDE, IOMMU controls, and trusted MSI as cooperating parts of a
confidential I/O design~\cite{riscv_cove_io_2026}. The important survey point
is that no single layer is sufficient. A measured device without constrained
DMA can still leak memory. A constrained DMA path without device identity can
still attach the wrong endpoint. A protected link without lifecycle control can
still carry traffic for an interface that was not securely assigned.

\subsubsection{Attested Service Channels}

SPDM and TDISP focus on the device and device-interface layers, but confidential
workloads also expose application-level endpoints. A key broker, model-serving
endpoint, storage proxy, packet-processing service, or DPU-hosted TLS
termination service may all need to prove that the network channel reaches the
expected protected code. \textbf{TLS+RA} integrates remote attestation into the
TLS handshake while keeping certificate-based endpoint identity and
attestation-based platform evidence independently deployable and independently
failing~\cite{weinhold2025tlsra}. This gives the survey a seventh layer:
attested service-channel binding. It is distinct from device identity because
the question is not only "which device is present" but "where does this
application connection terminate."

\subsection{Data Plane: DMA, Links, and Fabrics}

\subsubsection{IOMMU, SMMU, IOPMP, and Address Authority}

DMA mediation is the first enforcement layer for confidential I/O. Arm systems
use SMMU-based translation and protection to constrain device accesses in the
system memory path~\cite{arm_smmu_spec}. RISC-V systems use the IOMMU
specification for device address translation and IOPMP-style mechanisms for
requester-aware transaction filtering~\cite{riscv_iommu_2023,riscv_iopmp_2026}.
Academic work such as sIOPMP emphasizes the scalability pressure: protected
systems may have many devices, many domains, and frequent mapping changes, so
the I/O protection mechanism must be both precise and efficient~\cite{feng2024siopmp}.

These mechanisms should be classified as access-control and translation
mechanisms. They do not automatically encrypt data on a link, validate the
device's internal firmware, or prove the endpoint state to a remote verifier.
Their job is to ensure that a bus master can access only the memory and MMIO
resources assigned to it. This classification prevents a common survey error:
equating any I/O protection with a complete trusted I/O system.

\subsubsection{PCIe IDE, CXL, RDMA, and Remote Memory Paths}

Link and fabric protection answer a different question: what happens to data as
it traverses PCIe, CXL, RDMA, or related interconnects. PCIe IDE provides
integrity and data encryption for PCIe traffic~\cite{pcie_ide}. CXL changes the
memory hierarchy by allowing memory expansion, pooling, and tiering beyond
local DRAM~\cite{cxl_spec,gouk2022directcxl,zhong2024cxltiers}. RDMA and
remote paging systems move application memory over the network or fabric with
minimal CPU involvement~\cite{wang2025odrp}. These technologies are not all
confidential-computing defenses by themselves, but they define the paths that a
confidential-computing defense must cover.

The design consequence is that confidential workloads need a data-path view, not
only a CPU view. A page that is protected while resident in local memory may
later move through a CXL fabric, be staged through an RDMA-capable NIC, or be
referenced by a device queue. Each transition introduces a new enforcement
point. The platform must decide whether protection comes from link encryption,
DMA access control, endpoint TEE support, software copying through trusted
buffers, or a combination of these mechanisms.

Secure storage over network fabrics adds freshness and crash-consistency
questions to this path. \textbf{Hazel} studies NVMe-over-Fabrics storage under a
confidential-computing threat model and combines counter leasing, NVMe metadata,
a disaggregated Merkle-tree design, and BlueField-3 offload to provide
confidentiality, integrity, and freshness without changing the NVMe-oF
protocol~\cite{chrapek2026hazel}. This kind of work is useful for the survey
because it shows that a protected I/O path must sometimes defend not only
in-flight packets, but also remote persistent state and replay freshness.

\subsection{Trusted I/O Systems and Device TEEs}

\subsubsection{SEV-TIO and VM-Oriented Trusted I/O}

AMD SEV-TIO provides a concrete industry design point for trusted I/O in
encrypted-VM systems~\cite{amd_sev_tio_2023}. It is useful in this survey even
though the main architectural focus is Arm and RISC-V, because it clarifies the
general pattern: a confidential VM needs device assignment, device trust
establishment, protected DMA, interrupt handling, and attestation-compatible
state. These requirements reappear in CoVE-IO and in Arm CCA-oriented device
work.

The SEV-TIO comparison also helps avoid an overly narrow interpretation of
trusted I/O. Trusted I/O is not merely "attach a device to a VM." It is the
combination of device identity, interface lifecycle, memory access restriction,
link security, interrupt delivery, and evidence that a verifier can evaluate.
Different platforms divide those responsibilities differently, but the required
security properties are similar.

\subsubsection{Accelerator TEEs and Confidential Offload}

Accelerator TEEs push the same problem into device-local computation. The
accelerator may execute kernels, hold model weights, stage tensors, manage
queues, or perform packet-processing offload. A CPU confidential VM cannot
automatically protect that state once it leaves the CPU package. The accelerator
itself must either become part of the trusted boundary or be treated as
untrusted with protocols that keep secrets inside the confidential workload.

ACAI studies accelerator protection in the Arm CCA context~\cite{acai2023}.
PORTAL studies protected Arm CCA device access for mobile SoCs, using CCA
memory isolation to avoid repeated shared-buffer encryption~\cite{sang2025portal}.
CAGE adapts GPU and FPGA accelerator workflows to Arm CCA with shadow tasks and
GPC/GPT-backed checks~\cite{wang2026cage}. StrongBox is an earlier Arm GPU TEE
baseline that uses TrustZone and Stage-2 translation to protect unified-memory
GPU computation~\cite{deng2022strongbox}.
ITX demonstrates confidential computing within an AI accelerator, giving a
concrete example of device-local confidential execution~\cite{vaswani2023itx}.
The accelerator TEE SoK organizes this space around device identity, memory and
bus protection, runtime TCB, scheduling, and attestation~\cite{sok-tee}. These
works are important for network defense as well because modern NICs, DPUs, and
SmartNICs increasingly resemble accelerators: they execute firmware and
programmable code, maintain queues, and transform data before the CPU sees it.

\subsubsection{SmartNIC, DPU, and Rack-Scale Offload}

SmartNICs and DPUs sit between the device-interface layer and the application
service layer. They can expose SR-IOV functions and virtual switches to the
host, run embedded Arm-side control software, terminate tunnels, process
packets, or mediate storage traffic. HETEE and later cloud-scale work show one
route for protecting accelerator and DPU-like offload: place a security
controller on the data path so that non-TEE devices can be assigned, measured,
and cleaned up as part of a confidential workload lifecycle~\cite{zhu2020hetee,dhar2024cloudscale}.
Commercial DPU documentation shows another route: DPU-local roots such as
OP-TEE/fTPM and replay-protected storage can provide platform identity and key
protection building blocks~\cite{nvidia_bluefield_operation_2025}. That
commercial evidence should be used carefully, because the cited BlueField
documentation describes fTPM support and its limitations rather than a complete
production confidential networking TEE.

S-NIC and TNIC fill two additional gaps. S-NIC answers how a SmartNIC can
isolate tenant network functions from other functions and from the NIC OS, while
TNIC answers how a NIC can provide a small, verifiable trust root for networked
protocols~\cite{zhou2024snic,giantsidi2025tnic}. Hazel then shows how DPU/SmartNIC
offload can be used in a secure storage data path while preserving freshness
properties~\cite{chrapek2026hazel}. Together these papers make the DPU/NIC
category more precise: one must distinguish local resource isolation, endpoint
attestation, protocol-level non-equivocation, and storage freshness.

For survey classification, these systems are neither ordinary firewall/IDS
papers nor complete replacements for CoVE-IO or TDISP. They are evidence that
network offload devices are becoming programmable trust endpoints. A secure
confidential networking stack must therefore compose DPU-local roots,
device-interface lifecycle, DMA restrictions, link protection, and attested
application-channel termination.

\subsection{Taxonomy}

\begin{table*}[t]
\centering
\caption{Confidential I/O and Network Defense Taxonomy}
\label{tab:confidential_io_taxonomy}
\begin{tabular}{@{}p{0.18\linewidth}p{0.26\linewidth}p{0.25\linewidth}p{0.22\linewidth}@{}}
\toprule
\textbf{Layer} & \textbf{Primary Question} & \textbf{Representative Mechanisms} & \textbf{Survey Role} \\ \midrule
Device identity and measurement & Which endpoint is participating, and what state can it prove? & SPDM, SPDM secured messages, SPDM over TCP & Establish evidence and secure sessions for device trust. \\
Device-interface lifecycle & Which interface is assigned to the confidential workload? & TDISP, trusted device interface state, CoVE-IO TDI/TDM/DSM & Bind a device function or interface to a protected domain. \\
Attested service channel & Does the application channel terminate in the attested protected endpoint? & TLS+RA, RA-TLS-style endpoint binding, verifier policy & Bind network services, KBS endpoints, proxies, and DPU-hosted TLS termination to TEE evidence. \\
DMA and MMIO authority & Which memory and registers may the device access? & SMMU, RISC-V IOMMU, IOPMP, sIOPMP & Enforce address translation and transaction filtering. \\
Interrupt and event delivery & Can completion and error events be delivered without host spoofing or confusion? & AIA, trusted MSI, Devlore-style interrupt ownership & Keep device ownership coherent after asynchronous events. \\
Link and fabric protection & What protects traffic across PCIe, CXL, RDMA, and remote-memory paths? & PCIe IDE, CXL fabric controls, protected RDMA/remote paging designs & Protect data once it leaves CPU-local memory paths. \\
NIC-local root and function isolation & Can a SmartNIC/DPU isolate tenant functions or provide a trusted network primitive? & S-NIC, TNIC, DPU OP-TEE/fTPM & Treat NICs and DPUs as explicit trust endpoints rather than opaque offload boxes. \\
Confidential storage data path & Can remote storage preserve confidentiality, integrity, and freshness? & Hazel, protected NVMe-oF metadata, DPU crypto offload & Extend confidential I/O beyond packets into persistent remote state. \\
Device-local confidential execution & Does the accelerator or offload engine hold secrets or execute trusted code? & ACAI, PORTAL, CAGE, StrongBox, ITX, HETEE, CloudScale, BlueField OP-TEE/fTPM, accelerator TEE designs & Extend or avoid extending the TCB to accelerators and offload devices. \\ \bottomrule
\end{tabular}
\end{table*}

Table~\ref{tab:confidential_io_taxonomy} captures the classification used in
the rest of this survey. The table is intentionally layered because
confidential I/O failures are often composition failures. A system can have
SPDM but weak DMA protection, PCIe IDE but no trusted device-interface
lifecycle, or a device TEE but an ambiguous interrupt path. The practical
question is whether all relevant layers preserve the same ownership boundary.

\subsection{Summary}

Confidential I/O and network defense is the point where confidential computing
becomes a whole-platform problem. SPDM and secured messages provide identity,
measurement, and session foundations. TDISP and CoVE-IO introduce trusted
device-interface lifecycle concepts. TLS+RA-style protocols bind application
channels to attested endpoints. SMMU, IOMMU, IOPMP, and sIOPMP constrain DMA
authority. AIA and trusted MSI address interrupt delivery. PCIe IDE, CXL, RDMA,
NVMe-oF, and remote-memory systems define the fabric paths that protected data
may cross. Accelerator, SmartNIC, NIC, and DPU work such as ACAI, PORTAL, CAGE,
StrongBox, ITX, HETEE, CloudScale, S-NIC, TNIC, Hazel, and BlueField
OP-TEE/fTPM shows that device-local, NIC-local, storage-path, or
controller-mediated execution can become part of the confidential boundary. A
survey that focuses only on CPU TEE mechanisms would miss this entire class of
defenses; for modern Arm and RISC-V confidential-computing platforms, these
mechanisms are core rather than peripheral.
